\section{N}

% nähen (to sew)
\begin{verbentry}{
  style = {default},
  toc = {nähen (to sew)},
  name = {nähen},
  english = {to sew},
  type = {regular},
  auxiliary = {haben}
}


\vspace{5pt}
\hrule
\vspace{5pt}

\begin{tabular}{rl}
\multicolumn{2}{l}{\textbf{PRÄSENS}}\\
ich & \textbf{nähe}\\
du & \textbf{nähst}\\
er/sie/es & \textbf{näht}\\
wir & \textbf{nähen}\\
ihr & \textbf{näht}\\
sie/Sie & \textbf{nähen}\\
\end{tabular}
\hfill
\begin{tabular}{rl}
\multicolumn{2}{l}{\textbf{PRÄTERITUM}}\\
ich & \textbf{nähte}\\
du & \textbf{nähntest}\\
er/sie/es & \textbf{nähte}\\
wir & \textbf{nähten}\\
ihr & \textbf{nähtet}\\
sie/Sie & \textbf{nähten}\\
\end{tabular}
\hfill
\begin{tabular}{rl}
\multicolumn{2}{l}{\textbf{IMPERATIV}}\\
& \textbf{nähe (du)!}\\
& \textbf{nähen wir!}\\
& \textbf{näht (ihr)!}\\
& \textbf{nähen Sie!}\\
\end{tabular}

\vspace{10pt}

\begin{tabular}{rl}
\multicolumn{2}{l}{\textbf{PERFEKT}}\\
& haben + \textbf{genäht}\\
\end{tabular}
\hfill
\begin{tabular}{rl}
\multicolumn{2}{l}{\textbf{FUTURE I}}\\
& werden + \textbf{nähen}\\
\end{tabular}
\hfill
\begin{tabular}{rl}
\multicolumn{2}{l}{\textbf{PLUSQUAMPERFEKT}}\\
& hatten + \textbf{genäht}\\
\end{tabular}

\vspace{10pt}

\begin{tabular}{lll}
\multicolumn{3}{l}{\textbf{MIT PRÄPOSITIONEN}}\\
& \textbf{nähen an} + Dat & (to sew on / at)\\
& \textbf{nähen für} + Akk & (to sew for / on behalf of)\\
\end{tabular}
\hfill
\begin{tabular}{lll}
\multicolumn{3}{l}{\textbf{TRENNBARE VERBEN}}\\
& \textbf{auf$|$nähen} & (to sew onto / attach by sewing)\\
\end{tabular}
\vspace{-5pt}
\end{verbentry}

% nehmen (to take)
\begin{verbentry}{
  style = {acc},
  toc = {nehmen (to take)},
  name = {nehmen},
  english = {to take},
  type = {irregular, + Akkusative},
  auxiliary = {haben}
}
 
    
     \vspace{5pt}

    \hrule
    
    \vspace{5pt}

  \begin{tabular}{rl}
     \multicolumn{2}{l}{\textbf{PRÄSENS} }\\
    ich & \textbf{nehme}\\
    du & \textbf{nimmst}\\
    er/sie/es & \textbf{nimmt}\\
    wir & \textbf{nehmen}\\
    ihr & \textbf{nehmt}\\
    sie/Sie & \textbf{nehmen}\\
  \end{tabular}
  \hfill
     \begin{tabular}{rl}
    \multicolumn{2}{l}{\textbf{PRÄTERITUM} }\\
    ich & \textbf{nahm}\\
    du & \textbf{nahmst}\\
    er/sie/es & \textbf{nahm}\\
    wir & \textbf{nahmen}\\
    ihr & \textbf{nahmt}\\
    sie/Sie & \textbf{nahmen}\\
  \end{tabular}
  \hfill
  \begin{tabular}{rl}
     \multicolumn{2}{l}{\textbf{IMPERATIV} }\\
   & \textbf{nimm (du)!}\\
   & \textbf{nehmen wir!}\\
   & \textbf{nehmt (ihr)!}\\
   & \textbf{nehmen Sie!}\\
   & \vspace{10pt}
  \end{tabular}

    \vspace{10pt}

 \begin{tabular}{rl}
     \multicolumn{2}{l}{\textbf{PERFEKT}}\\
  &  haben + \textbf{genommen}\\
  \end{tabular}
\hfill
   \begin{tabular}{rl}
   
    \multicolumn{2}{l}{\textbf{FUTURE I} }\\
  &  werden + \textbf{nehmen}\\
 
  \end{tabular}
\hfill
   \begin{tabular}{rl}
      \multicolumn{2}{l}{\textbf{PLUSQUAMPERFEKT}}\\
  &  hatten + \textbf{genommen}\\
  \end{tabular}

    \vspace{10pt}

   \begin{tabular}{lll}
      \multicolumn{3}{l}{\textbf{MIT PRÄPOSITIONEN}}\\
& \textbf{nehmen aus} + Dat & (to take from)\\
& \textbf{nehmen von} + Dat & (to take from)\\
&\textbf{nehmen für} + Akk & (to take for)\\
&\textbf{nehmen über} + Akk & (to take over / control of)\\
\vspace{50pt}
  \end{tabular}
  \hfill
   \begin{tabular}{lll}
      \multicolumn{3}{l}{\textbf{TRENNBARE VERBEN}}\\
      & \textbf{ab$|$nehmen} & (to turn off / lose weight)\\
 & \textbf{an$|$nehmen} & (to accept / assume)\\
 & \textbf{auf$|$nehmen} & (to pick up / record)\\
 & \textbf{ein$|$nehmen} & (to take something / to occupy)\\
  & \textbf{mit$|$nehmen} & (to take along)\\
  & \textbf{über$|$nehmen} & (to take over)\\
  & \textbf{weg$|$nehmen} & (to take away)\\
     & \textbf{zurück$|$nehmen} & (to take back)\\
     & \textbf{teil$|$nehmen} & (to participate)\\
       & \textbf{zu$|$nehmen} & (to gain weight)\\
  \end{tabular}



  \vspace{-5pt}
\end{verbentry}

% nennen (to name)
\begin{verbentry}{
  style = {acc},
  toc = {nennen (to name)},
  name = {nennen},
  english = {to name},
  type = {irregular, + Akkusativ},
  auxiliary = {haben}
}
 
    
     \vspace{5pt}

    \hrule
    
    \vspace{5pt}

  \begin{tabular}{rl}
     \multicolumn{2}{l}{\textbf{PRÄSENS} }\\
    ich & \textbf{nenne}\\
    du & \textbf{nennst}\\
    er/sie/es & \textbf{nennt}\\
    wir & \textbf{nennen}\\
    ihr & \textbf{nennt}\\
    sie/Sie & \textbf{nennen}\\
  \end{tabular}
  \hfill
     \begin{tabular}{rl}
    \multicolumn{2}{l}{\textbf{PRÄTERITUM} }\\
    ich & \textbf{nannte}\\
    du & \textbf{nanntest}\\
    er/sie/es & \textbf{nannte}\\
    wir & \textbf{nannten}\\
    ihr & \textbf{nanntet}\\
    sie/Sie & \textbf{nannten}\\
  \end{tabular}
  \hfill
  \begin{tabular}{rl}
     \multicolumn{2}{l}{\textbf{IMPERATIV} }\\
   & \textbf{nenn(e) (du)!}\\
   & \textbf{nennen wir!}\\
   & \textbf{nennt (ihr)!}\\
   & \textbf{nennen Sie!}\\
   & \vspace{10pt}
  \end{tabular}

    \vspace{10pt}

 \begin{tabular}{rl}
     \multicolumn{2}{l}{\textbf{PERFEKT}}\\
  &  haben + \textbf{genannt}\\
  \end{tabular}
\hfill
   \begin{tabular}{rl}
   
    \multicolumn{2}{l}{\textbf{FUTURE I} }\\
   & werden + \textbf{nennen}\\
 
  \end{tabular}
\hfill
   \begin{tabular}{rl}
      \multicolumn{2}{l}{\textbf{PLUSQUAMPERFEKT}}\\
   & hatten + \textbf{genannt}\\
  \end{tabular}
  
  
\vspace{10pt}

\begin{tabular}{lll}
\multicolumn{3}{l}{\textbf{MIT PRÄPOSITIONEN}}\\
& \textbf{---} & \\
\end{tabular}
\hfill
\begin{tabular}{lll}
\multicolumn{3}{l}{\textbf{TRENNBARE VERBEN}}\\
& \textbf{---} & \\
\end{tabular}
\vspace{-5pt}
\end{verbentry}

% niesen (to sneeze)
\begin{verbentry}{
  style = {default},
  toc = {niesen (to sneeze)},
  name = {niesen},
  english = {to sneeze},
  type = {regular},
  auxiliary = {haben}
}


\vspace{5pt}
\hrule
\vspace{5pt}

\begin{tabular}{rl}
\multicolumn{2}{l}{\textbf{PRÄSENS}}\\
ich & \textbf{niese}\\
du & \textbf{niest}\\
er/sie/es & \textbf{niest}\\
wir & \textbf{niesen}\\
ihr & \textbf{niest}\\
sie/Sie & \textbf{niesen}\\
\end{tabular}
\hfill
\begin{tabular}{rl}
\multicolumn{2}{l}{\textbf{PRÄTERITUM}}\\
ich & \textbf{nieste}\\
du & \textbf{niestest}\\
er/sie/es & \textbf{nieste}\\
wir & \textbf{niesten}\\
ihr & \textbf{niestet}\\
sie/Sie & \textbf{niesten}\\
\end{tabular}
\hfill
\begin{tabular}{rl}
\multicolumn{2}{l}{\textbf{IMPERATIV}}\\
& \textbf{niese (du)!}\\
& \textbf{niesen wir!}\\
& \textbf{niest (ihr)!}\\
& \textbf{niesen Sie!}\\
\end{tabular}

\vspace{10pt}

\begin{tabular}{rl}
\multicolumn{2}{l}{\textbf{PERFEKT}}\\
& haben + \textbf{geniest}\\
\end{tabular}
\hfill
\begin{tabular}{rl}
\multicolumn{2}{l}{\textbf{FUTURE I}}\\
& werden + \textbf{niesen}\\
\end{tabular}
\hfill
\begin{tabular}{rl}
\multicolumn{2}{l}{\textbf{PLUSQUAMPERFEKT}}\\
& hatten + \textbf{geniest}\\
\end{tabular}

\vspace{10pt}

\begin{tabular}{lll}
\multicolumn{3}{l}{\textbf{MIT PRÄPOSITIONEN}}\\
& \textbf{niesen in} + Akk & (to sneeze into)\\
\end{tabular}
\hfill
\begin{tabular}{lll}
\multicolumn{3}{l}{\textbf{TRENNBARE VERBEN}}\\
& \textbf{---} & \\
\end{tabular}
\vspace{-5pt}
\end{verbentry}

% nutzen (to use)
\begin{verbentry}{
  style = {acc},
  toc = {nutzen (to use)},
  name = {nutzen},
  english = {to use},
  type = {regular, + Akkusativ},
  auxiliary = {haben}
}
 
    
     \vspace{5pt}

    \hrule
    
    \vspace{5pt}

  \begin{tabular}{rl}
     \multicolumn{2}{l}{\textbf{PRÄSENS} }\\
    ich & \textbf{nutze}\\
    du & \textbf{nutzt}\\
    er/sie/es & \textbf{nutzt}\\
    wir & \textbf{nutzen}\\
    ihr & \textbf{nutzt}\\
    sie/Sie & \textbf{nutzen}\\
  \end{tabular}
  \hfill
     \begin{tabular}{rl}
    \multicolumn{2}{l}{\textbf{PRÄTERITUM} }\\
    ich & \textbf{nutzte}\\
    du & \textbf{nutztest}\\
    er/sie/es & \textbf{nutzte}\\
    wir & \textbf{nutzten}\\
    ihr & \textbf{nutztet}\\
    sie/Sie & \textbf{nutzten}\\
  \end{tabular}
  \hfill
  \begin{tabular}{rl}
     \multicolumn{2}{l}{\textbf{IMPERATIV} }\\
   & \textbf{nutz(e) (du)!}\\
   & \textbf{ven wir!}\\
   & \textbf{nutzt (ihr)!}\\
   & \textbf{nutzen Sie!}\\
   & \vspace{10pt}
  \end{tabular}

    \vspace{10pt}

 \begin{tabular}{rl}
     \multicolumn{2}{l}{\textbf{PERFEKT}}\\
  &  haben + \textbf{genutzt}\\
  \end{tabular}
\hfill
   \begin{tabular}{rl}
   
    \multicolumn{2}{l}{\textbf{FUTURE I} }\\
  &  werden + \textbf{nutzen}\\
 
  \end{tabular}
\hfill
   \begin{tabular}{rl}
      \multicolumn{2}{l}{\textbf{PLUSQUAMPERFEKT}}\\
  &  hatten + \textbf{genutzt}\\
  \end{tabular}
  
  
\vspace{10pt}

\begin{tabular}{lll}
\multicolumn{3}{l}{\textbf{MIT PRÄPOSITIONEN}}\\
& \textbf{---} & \\
\end{tabular}
\hfill
\begin{tabular}{lll}
\multicolumn{3}{l}{\textbf{TRENNBARE VERBEN}}\\
& \textbf{---} & \\
\end{tabular}
\vspace{-5pt}
\end{verbentry}
